\documentclass{beamer}
\usetheme{Madrid}
\usecolortheme{default}

% Define custom colors for AI slides
\definecolor{AIBlue}{RGB}{30, 144, 255}
\definecolor{AILightBlue}{RGB}{173, 216, 230}
\definecolor{AINavy}{RGB}{25, 25, 112}

% Commands to switch to AI color scheme
\newcommand{\AIcolors}{%
  \setbeamercolor{frametitle}{bg=AIBlue,fg=white}%
  \setbeamercolor{title}{bg=AINavy,fg=white}%
  \setbeamercolor{structure}{fg=AIBlue}%
  \setbeamercolor{item}{fg=AINavy}%
  \setbeamercolor{normal text}{fg=black}%
}

% Command to restore default colors
\newcommand{\defaultcolors}{%
  \setbeamercolor{frametitle}{bg=osaorange,fg=white}%
  \setbeamercolor{title}{bg=osaorange,fg=white}%
  \setbeamercolor{structure}{fg=osaorange}%
  \setbeamercolor{item}{fg=black}%
  \setbeamercolor{normal text}{fg=black}%
}

% Define OSU orange color
\definecolor{osaorange}{RGB}{220,115,38}

% Set colors
\setbeamercolor{structure}{fg=osaorange}
\setbeamercolor{title}{fg=white,bg=osaorange}
\setbeamercolor{frametitle}{fg=white,bg=osaorange}

\usepackage{graphicx}
\usepackage{amsmath}
\usepackage{booktabs}
\usepackage{multirow}
\usepackage{hyperref}
\usepackage{tikz}
\usetikzlibrary{shapes,arrows}

\title{H524 Introduction to Biostatistics}
\subtitle{Week 5: Hypothesis Testing}
\author{Dr. John Molitor}
\institute{}
\date{Fall 2025}

\begin{document}

\frame{\titlepage}

\begin{frame}
\frametitle{Outline}
\tableofcontents
\end{frame}

\section{Review from Week 4}

\begin{frame}
\frametitle{Quick Review: Confidence Intervals}
\textbf{Last week we covered:}
\begin{itemize}
\item Point estimates vs interval estimates
\item Confidence intervals for population mean
\item General form: $\bar{x} \pm t_{\alpha/2, n-1} \times \frac{s}{\sqrt{n}}$
\item Interpretation: 95\% of intervals will contain $\mu$
\item Trade-off between confidence level and precision
\end{itemize}

\vspace{0.5cm}
\textbf{This Week:} Hypothesis testing - making decisions about population parameters

\vspace{0.3cm}
\textbf{Reading:} Pagano \& Gauvreau Chapter 10
\end{frame}

\begin{frame}
\frametitle{From Estimation to Testing}
\textbf{Two approaches to statistical inference:}

\vspace{0.5cm}
\textbf{1. Estimation (Weeks 3-4):}
\begin{itemize}
\item What is the value of the parameter?
\item Construct confidence interval
\item Example: ``True mean cholesterol is between 195 and 205 mg/dL''
\end{itemize}

\vspace{0.5cm}
\textbf{2. Hypothesis Testing (This week):}
\begin{itemize}
\item Is a specific claim about the parameter true?
\item Make a decision: reject or fail to reject
\item Example: ``Is the true mean cholesterol greater than 200 mg/dL?''
\end{itemize}

\vspace{0.5cm}
\textbf{Both use:} Sample data, sampling distributions, probability theory
\end{frame}

\section{Introduction to Hypothesis Testing (Ch 10.1)}

\begin{frame}
\frametitle{What is a Hypothesis Test?}
\textbf{Definition:} A statistical procedure to assess evidence about a population parameter

\vspace{0.5cm}
\textbf{Basic idea:}
\begin{enumerate}
\item Make a claim about a population parameter
\item Collect sample data
\item Use probability to decide if data supports the claim
\item Make a decision: reject or fail to reject the claim
\end{enumerate}

\vspace{0.5cm}
\textbf{Medical example:}
\begin{itemize}
\item Claim: ``New drug lowers blood pressure by more than 10 mmHg on average''
\item Collect data from clinical trial
\item Calculate probability of observing this data if drug has no effect
\item If probability is very small, reject ``no effect'' claim
\end{itemize}
\end{frame}

\begin{frame}
\frametitle{The Logic of Hypothesis Testing}
\textbf{Key principle:} \textit{Proof by contradiction}

\vspace{0.5cm}
\textbf{Analogy - Criminal trial:}
\begin{itemize}
\item \textbf{Null hypothesis:} Defendant is innocent (default assumption)
\item \textbf{Alternative hypothesis:} Defendant is guilty (what we're trying to prove)
\item \textbf{Evidence:} Trial data
\item \textbf{Decision:} Guilty (reject null) or Not Guilty (fail to reject null)
\item \textbf{Standard:} Beyond reasonable doubt (low probability threshold)
\end{itemize}

\vspace{0.5cm}
\textbf{Important:} ``Not guilty'' $\neq$ ``innocent''

Similarly: ``Fail to reject $H_0$'' $\neq$ ``$H_0$ is true''

\vspace{0.3cm}
We never ``accept'' the null hypothesis - we only fail to reject it!
\end{frame}

\begin{frame}
\frametitle{Null and Alternative Hypotheses}
\textbf{Null Hypothesis ($H_0$):}
\begin{itemize}
\item Statement of ``no effect'' or ``no difference''
\item Status quo, default assumption
\item What we assume is true unless proven otherwise
\item Always contains equality: $=$, $\leq$, or $\geq$
\end{itemize}

\vspace{0.5cm}
\textbf{Alternative Hypothesis ($H_A$ or $H_1$):}
\begin{itemize}
\item Statement of ``effect exists'' or ``difference exists''
\item What the researcher wants to demonstrate
\item What we conclude if evidence is strong enough
\item Contains inequality: $\neq$, $>$, or $<$
\end{itemize}

\vspace{0.5cm}
\textbf{Key rule:} $H_0$ and $H_A$ must be:
\begin{itemize}
\item Mutually exclusive (can't both be true)
\item Exhaustive (one must be true)
\end{itemize}
\end{frame}

\begin{frame}
\frametitle{Types of Hypotheses: Two-Sided vs One-Sided}
\textbf{Two-sided (two-tailed) test:}
$$H_0: \mu = \mu_0 \quad \text{vs} \quad H_A: \mu \neq \mu_0$$

\begin{itemize}
\item Testing if parameter differs from null value (either direction)
\item Example: ``Is mean cholesterol different from 200 mg/dL?''
\item Most common, more conservative
\end{itemize}

\vspace{0.5cm}
\textbf{One-sided (one-tailed) test - Upper tail:}
$$H_0: \mu \leq \mu_0 \quad \text{vs} \quad H_A: \mu > \mu_0$$

\begin{itemize}
\item Testing if parameter is greater than null value
\item Example: ``Is mean blood pressure higher than 120 mmHg?''
\end{itemize}

\vspace{0.5cm}
\textbf{One-sided test - Lower tail:}
$$H_0: \mu \geq \mu_0 \quad \text{vs} \quad H_A: \mu < \mu_0$$

\begin{itemize}
\item Testing if parameter is less than null value
\item Example: ``Does drug reduce cholesterol below 200 mg/dL?''
\end{itemize}
\end{frame}

\begin{frame}
\frametitle{Example: Setting Up Hypotheses}
\textbf{Scenario 1:} Normal body temperature is claimed to be 98.6°F. We want to test if this is accurate.

\vspace{0.3cm}
$$H_0: \mu = 98.6 \quad \text{vs} \quad H_A: \mu \neq 98.6$$

\vspace{0.5cm}
\textbf{Scenario 2:} A new drug claims to lower blood pressure. Current mean is 140 mmHg.

\vspace{0.3cm}
$$H_0: \mu \geq 140 \quad \text{vs} \quad H_A: \mu < 140$$

\vspace{0.5cm}
\textbf{Scenario 3:} We suspect exposure to a toxin increases cancer rate above baseline 5\%.

\vspace{0.3cm}
$$H_0: p \leq 0.05 \quad \text{vs} \quad H_A: p > 0.05$$

\vspace{0.5cm}
\textbf{Key decision:} Choose one-sided vs two-sided \textit{before} collecting data, based on research question!
\end{frame}

\section{Test Statistics and P-values (Ch 10.2)}

\begin{frame}
\frametitle{The Test Statistic}
\textbf{Definition:} A value calculated from sample data that measures how far the sample is from the null hypothesis

\vspace{0.5cm}
\textbf{General form:}
$$\text{Test Statistic} = \frac{\text{Estimate} - \text{Null Value}}{\text{Standard Error}}$$

\vspace{0.5cm}
\textbf{For testing mean (large sample or $\sigma$ known):}
$$Z = \frac{\bar{x} - \mu_0}{\sigma/\sqrt{n}}$$

\vspace{0.5cm}
\textbf{For testing mean ($\sigma$ unknown, small sample):}
$$t = \frac{\bar{x} - \mu_0}{s/\sqrt{n}}$$

\vspace{0.3cm}
where $\mu_0$ is the null hypothesis value
\end{frame}

\begin{frame}
\frametitle{Interpreting the Test Statistic}
\textbf{What does the test statistic tell us?}

\vspace{0.3cm}
\begin{itemize}
\item Measures ``how many standard errors away'' the sample mean is from null value
\item Large $|t|$ or $|Z|$: Sample mean far from null value
\item Small $|t|$ or $|Z|$: Sample mean close to null value
\end{itemize}

\vspace{0.5cm}
\textbf{Examples:}
\begin{itemize}
\item $t = 0.5$: Sample mean is 0.5 SE from null (not unusual)
\item $t = 2.5$: Sample mean is 2.5 SE from null (unusual!)
\item $t = -3.8$: Sample mean is 3.8 SE below null (very unusual!)
\end{itemize}

\vspace{0.5cm}
\textbf{Question:} How large does $|t|$ need to be for us to reject $H_0$?

\vspace{0.3cm}
\textbf{Answer:} Depends on:
\begin{itemize}
\item Degrees of freedom (sample size)
\item Significance level (how much risk we'll accept)
\item One-sided vs two-sided test
\end{itemize}
\end{frame}

\begin{frame}
\frametitle{The P-value}
\textbf{Definition:} The probability of observing a test statistic as extreme as (or more extreme than) what we observed, \textit{assuming $H_0$ is true}

\vspace{0.5cm}
\textbf{For two-sided test:}
$$p\text{-value} = \Pr(|T| \geq |t_{obs}| \mid H_0 \text{ true})$$

\vspace{0.3cm}
\textbf{For one-sided test (upper tail):}
$$p\text{-value} = \Pr(T \geq t_{obs} \mid H_0 \text{ true})$$

\vspace{0.3cm}
\textbf{For one-sided test (lower tail):}
$$p\text{-value} = \Pr(T \leq t_{obs} \mid H_0 \text{ true})$$

\vspace{0.5cm}
\textbf{Interpretation:}
\begin{itemize}
\item Small p-value: Observed data would be very unlikely if $H_0$ were true
\item Large p-value: Observed data consistent with $H_0$
\end{itemize}
\end{frame}

\begin{frame}
\frametitle{Understanding P-values}
\textbf{Common misconceptions:}

\vspace{0.3cm}
\begin{itemize}
\item $\times$ p-value is NOT $\Pr(H_0 \text{ is true})$
\item $\times$ p-value is NOT $\Pr(\text{data due to chance})$
\item $\times$ 1 - p-value is NOT $\Pr(H_A \text{ is true})$
\end{itemize}

\vspace{0.5cm}
\textbf{Correct interpretation:}

\vspace{0.3cm}
``If the null hypothesis were true, the probability of obtaining a test statistic as extreme as or more extreme than the one we observed is [p-value]''

\vspace{0.5cm}
\textbf{Rule of thumb for interpretation:}
\begin{itemize}
\item $p < 0.01$: Very strong evidence against $H_0$
\item $0.01 \leq p < 0.05$: Strong evidence against $H_0$
\item $0.05 \leq p < 0.10$: Weak evidence against $H_0$
\item $p \geq 0.10$: Little or no evidence against $H_0$
\end{itemize}
\end{frame}

\begin{frame}
\frametitle{Significance Level ($\alpha$)}
\textbf{Definition:} The probability threshold for rejecting $H_0$

\vspace{0.5cm}
\textbf{Decision rule:}
\begin{itemize}
\item If $p\text{-value} < \alpha$: Reject $H_0$ (statistically significant)
\item If $p\text{-value} \geq \alpha$: Fail to reject $H_0$ (not statistically significant)
\end{itemize}

\vspace{0.5cm}
\textbf{Common choices:}
\begin{itemize}
\item $\alpha = 0.05$ (most common in biostatistics)
\item $\alpha = 0.01$ (more stringent, reduces false positives)
\item $\alpha = 0.10$ (more lenient, exploratory studies)
\end{itemize}

\vspace{0.5cm}
\textbf{Important:} Choose $\alpha$ \textit{before} collecting data!

\vspace{0.3cm}
Changing $\alpha$ after seeing results is called ``p-hacking'' and invalidates the test
\end{frame}

\section{Type I and Type II Errors (Ch 10.3)}

\begin{frame}
\frametitle{Two Types of Errors}
\textbf{Hypothesis testing is a decision with uncertainty}

\vspace{0.5cm}
\begin{center}
\begin{tabular}{cc|cc}
\toprule
& & \multicolumn{2}{c}{\textbf{Reality (Unknown)}} \\
& & $H_0$ True & $H_0$ False \\
\midrule
\multirow{2}{*}{\textbf{Decision}} & Reject $H_0$ & \textcolor{red}{Type I Error} & \textcolor{blue}{Correct} \\
& & \textcolor{red}{($\alpha$)} & \textcolor{blue}{(Power)} \\
& Fail to Reject & \textcolor{blue}{Correct} & \textcolor{red}{Type II Error} \\
& & \textcolor{blue}{($1-\alpha$)} & \textcolor{red}{($\beta$)} \\
\bottomrule
\end{tabular}
\end{center}

\vspace{0.5cm}
\textbf{Type I Error:} Reject $H_0$ when it's actually true (false positive)
\begin{itemize}
\item Probability = $\alpha$ (significance level)
\item Example: Conclude drug works when it doesn't
\end{itemize}

\vspace{0.3cm}
\textbf{Type II Error:} Fail to reject $H_0$ when it's actually false (false negative)
\begin{itemize}
\item Probability = $\beta$
\item Example: Miss a real drug effect
\end{itemize}
\end{frame}

\begin{frame}
\frametitle{Type I and Type II Errors: Medical Examples}
\textbf{Example 1: Disease screening test}

\vspace{0.3cm}
$H_0$: Patient does not have disease\\
$H_A$: Patient has disease

\vspace{0.3cm}
\begin{itemize}
\item \textbf{Type I Error:} Diagnose disease when patient is healthy (false positive)
  \begin{itemize}
  \item Consequence: Unnecessary treatment, anxiety
  \end{itemize}
\item \textbf{Type II Error:} Miss disease when patient is sick (false negative)
  \begin{itemize}
  \item Consequence: Delayed treatment, disease progression
  \end{itemize}
\end{itemize}

\vspace{0.5cm}
\textbf{Example 2: Drug approval}

\vspace{0.3cm}
$H_0$: Drug has no effect\\
$H_A$: Drug is effective

\vspace{0.3cm}
\begin{itemize}
\item \textbf{Type I Error:} Approve ineffective drug
  \begin{itemize}
  \item Consequence: Patients receive useless treatment
  \end{itemize}
\item \textbf{Type II Error:} Reject effective drug
  \begin{itemize}
  \item Consequence: Beneficial treatment unavailable
  \end{itemize}
\end{itemize}
\end{frame}

\begin{frame}
\frametitle{Controlling Type I Error}
\textbf{Type I Error rate = $\alpha$ (significance level)}

\vspace{0.5cm}
We control Type I error by choosing $\alpha$:
\begin{itemize}
\item Set $\alpha = 0.05$: 5\% chance of false positive
\item Set $\alpha = 0.01$: 1\% chance of false positive (more conservative)
\item Set $\alpha = 0.10$: 10\% chance of false positive (more lenient)
\end{itemize}

\vspace{0.5cm}
\textbf{Trade-off:}
\begin{itemize}
\item Decreasing $\alpha$ (stricter) reduces Type I errors
\item But increases Type II errors ($\beta$) - harder to detect real effects
\item Can't make both arbitrarily small without increasing sample size
\end{itemize}

\vspace{0.5cm}
\textbf{Convention in biostatistics:}
\begin{itemize}
\item Typically prioritize controlling Type I error
\item $\alpha = 0.05$ is standard (sometimes 0.01 for critical decisions)
\item Better to miss a real effect than falsely claim one exists
\end{itemize}
\end{frame}

\begin{frame}
\frametitle{Statistical Power}
\textbf{Power} = Probability of correctly rejecting $H_0$ when $H_A$ is true

$$\text{Power} = 1 - \beta$$

\vspace{0.5cm}
\textbf{What affects power?}
\begin{enumerate}
\item \textbf{Sample size ($n$):} Larger $n$ $\Rightarrow$ higher power
\item \textbf{Effect size:} Larger true difference $\Rightarrow$ higher power
\item \textbf{Significance level ($\alpha$):} Larger $\alpha$ $\Rightarrow$ higher power
\item \textbf{Variability ($\sigma$):} Lower $\sigma$ $\Rightarrow$ higher power
\end{enumerate}

\vspace{0.5cm}
\textbf{Typical goal:} Power $\geq 0.80$ (80\%)

\vspace{0.3cm}
Means 80\% chance of detecting the effect if it truly exists

\vspace{0.5cm}
\textbf{Power analysis:} Calculate required sample size to achieve desired power
\begin{itemize}
\item Done \textit{before} collecting data
\item Ensures study is adequately sized to detect meaningful effects
\end{itemize}
\end{frame}

\section{One-Sample t-Test (Ch 10.4)}

\begin{frame}
\frametitle{One-Sample t-Test}
\textbf{Used when:}
\begin{itemize}
\item Testing hypothesis about a population mean $\mu$
\item Population SD ($\sigma$) is unknown
\item Sample is random and approximately normal (or $n$ large by CLT)
\end{itemize}

\vspace{0.5cm}
\textbf{Hypotheses:}

Two-sided: $H_0: \mu = \mu_0$ vs $H_A: \mu \neq \mu_0$

One-sided: $H_0: \mu \leq \mu_0$ vs $H_A: \mu > \mu_0$ (or opposite)

\vspace{0.5cm}
\textbf{Test statistic:}
$$t = \frac{\bar{x} - \mu_0}{s/\sqrt{n}}$$

\vspace{0.3cm}
Under $H_0$: $t \sim t_{n-1}$ (t-distribution with $n-1$ degrees of freedom)

\vspace{0.5cm}
\textbf{Decision:}
\begin{itemize}
\item Calculate p-value from t-distribution
\item Reject $H_0$ if $p < \alpha$
\end{itemize}
\end{frame}

\begin{frame}
\frametitle{Example: Body Temperature}
\textbf{Question:} Is mean body temperature really 98.6°F?

\vspace{0.3cm}
\textbf{Data:} Sample of $n = 130$ healthy adults
\begin{itemize}
\item Sample mean: $\bar{x} = 98.25$°F
\item Sample SD: $s = 0.73$°F
\end{itemize}

\vspace{0.5cm}
\textbf{Step 1: Set up hypotheses}
$$H_0: \mu = 98.6 \quad \text{vs} \quad H_A: \mu \neq 98.6$$

\vspace{0.3cm}
\textbf{Step 2: Calculate test statistic}
$$t = \frac{98.25 - 98.6}{0.73/\sqrt{130}} = \frac{-0.35}{0.064} = -5.47$$

\vspace{0.3cm}
\textbf{Step 3: Find p-value}

From t-distribution with df = 129: $p < 0.0001$ (two-sided)

\vspace{0.3cm}
\textbf{Step 4: Conclusion}

$p < 0.05$, so reject $H_0$. Strong evidence that mean body temperature is not 98.6°F.
\end{frame}

\begin{frame}
\frametitle{Example: Blood Pressure Drug}
\textbf{Question:} Does new drug lower systolic BP below 140 mmHg?

\vspace{0.3cm}
\textbf{Data:} $n = 25$ patients after treatment
\begin{itemize}
\item Sample mean: $\bar{x} = 135$ mmHg
\item Sample SD: $s = 12$ mmHg
\end{itemize}

\vspace{0.5cm}
\textbf{Step 1: Set up hypotheses (one-sided)}
$$H_0: \mu \geq 140 \quad \text{vs} \quad H_A: \mu < 140$$

\vspace{0.3cm}
\textbf{Step 2: Calculate test statistic}
$$t = \frac{135 - 140}{12/\sqrt{25}} = \frac{-5}{2.4} = -2.08$$

\vspace{0.3cm}
\textbf{Step 3: Find p-value (one-sided, lower tail)}

From $t_{24}$ distribution: $p = 0.024$

\vspace{0.3cm}
\textbf{Step 4: Conclusion}

$p = 0.024 < 0.05$, reject $H_0$. Evidence that drug lowers BP below 140 mmHg.
\end{frame}

\begin{frame}
\frametitle{Critical Values vs P-values}
\textbf{Two equivalent approaches to hypothesis testing:}

\vspace{0.5cm}
\textbf{Method 1: P-value approach} (modern, preferred)
\begin{enumerate}
\item Calculate test statistic $t$
\item Calculate p-value
\item Compare p-value to $\alpha$
\item Reject $H_0$ if $p < \alpha$
\end{enumerate}

\vspace{0.5cm}
\textbf{Method 2: Critical value approach} (classical)
\begin{enumerate}
\item Calculate test statistic $t$
\item Find critical value(s) from t-table for given $\alpha$
\item Compare $|t|$ to critical value
\item Reject $H_0$ if $|t| > t_{critical}$ (two-sided)
\end{enumerate}

\vspace{0.5cm}
\textbf{Advantage of p-values:}
\begin{itemize}
\item Provides exact strength of evidence
\item Can compare to any $\alpha$ level
\item Software gives p-values automatically
\end{itemize}
\end{frame}

\begin{frame}
\frametitle{Critical Values for Two-Sided Tests}
\textbf{For two-sided test at $\alpha = 0.05$:}

\vspace{0.3cm}
Reject $H_0$ if $|t| > t_{0.975, n-1}$

\vspace{0.5cm}
\begin{center}
\begin{tabular}{ccc}
\toprule
df & $t_{0.975}$ & Interpretation \\
\midrule
5 & 2.571 & $|t| > 2.571$ is significant \\
10 & 2.228 & $|t| > 2.228$ is significant \\
20 & 2.086 & $|t| > 2.086$ is significant \\
30 & 2.042 & $|t| > 2.042$ is significant \\
100 & 1.984 & $|t| > 1.984$ is significant \\
$\infty$ & 1.960 & $|t| > 1.960$ is significant \\
\bottomrule
\end{tabular}
\end{center}

\vspace{0.5cm}
\textbf{Note:} As df increases, $t$ critical values approach $Z = 1.96$

\vspace{0.3cm}
For large samples ($n > 30$), can use normal approximation
\end{frame}

\section{One-Sample z-Test (Ch 10.5)}

\begin{frame}
\frametitle{One-Sample z-Test}
\textbf{Used when:}
\begin{itemize}
\item Testing hypothesis about population mean $\mu$
\item Population SD ($\sigma$) is \textbf{known} (rare in practice!)
\item OR sample size is very large ($n > 100$)
\end{itemize}

\vspace{0.5cm}
\textbf{Test statistic:}
$$Z = \frac{\bar{x} - \mu_0}{\sigma/\sqrt{n}}$$

\vspace{0.3cm}
Under $H_0$: $Z \sim N(0, 1)$ (standard normal)

\vspace{0.5cm}
\textbf{Critical values for $\alpha = 0.05$:}
\begin{itemize}
\item Two-sided: $|Z| > 1.96$
\item One-sided: $Z > 1.645$ (upper) or $Z < -1.645$ (lower)
\end{itemize}

\vspace{0.5cm}
\textbf{In practice:} Almost always use t-test instead
\begin{itemize}
\item Rarely know $\sigma$ without knowing $\mu$
\item For large $n$, t-test and z-test give nearly identical results
\end{itemize}
\end{frame}

\begin{frame}
\frametitle{Example: IQ Test with Known SD}
\textbf{Scenario:} IQ tests are designed to have $\sigma = 15$ (known). We want to test if a special program increases mean IQ above 100.

\vspace{0.3cm}
\textbf{Data:} $n = 64$ students in program
\begin{itemize}
\item Sample mean: $\bar{x} = 103$
\item Known SD: $\sigma = 15$
\end{itemize}

\vspace{0.5cm}
\textbf{Hypotheses:}
$$H_0: \mu \leq 100 \quad \text{vs} \quad H_A: \mu > 100$$

\vspace{0.3cm}
\textbf{Test statistic:}
$$Z = \frac{103 - 100}{15/\sqrt{64}} = \frac{3}{1.875} = 1.60$$

\vspace{0.3cm}
\textbf{P-value:} $\Pr(Z > 1.60) = 0.055$

\vspace{0.3cm}
\textbf{Conclusion:} $p = 0.055 > 0.05$, fail to reject $H_0$. Marginal evidence that program increases IQ.
\end{frame}

\section{Confidence Intervals and Hypothesis Tests}

\begin{frame}
\frametitle{Relationship Between CIs and Hypothesis Tests}
\textbf{Key insight:} Confidence intervals and hypothesis tests are closely related!

\vspace{0.5cm}
\textbf{For two-sided test at level $\alpha$:}

\vspace{0.3cm}
Reject $H_0: \mu = \mu_0$ if and only if the $(1-\alpha) \times 100$\% CI does NOT contain $\mu_0$

\vspace{0.5cm}
\textbf{Example:} Blood pressure drug data
\begin{itemize}
\item $n = 25$, $\bar{x} = 135$ mmHg, $s = 12$ mmHg
\item 95\% CI: $135 \pm 2.064 \times \frac{12}{\sqrt{25}} = (130.0, 140.0)$
\item Test $H_0: \mu = 140$ vs $H_A: \mu \neq 140$
\item Since 140 is \textit{in} the CI, we would \textit{fail to reject} $H_0$ (for two-sided test)
\end{itemize}

\vspace{0.3cm}
\textbf{Note:} This equivalence only holds for two-sided tests!

\vspace{0.3cm}
For one-sided tests, use one-sided confidence bounds
\end{frame}

\begin{frame}
\frametitle{Using CIs for Multiple Comparisons}
\textbf{Advantage of confidence intervals:}

\vspace{0.3cm}
Can simultaneously test multiple null values!

\vspace{0.5cm}
\textbf{Example:} Cholesterol drug trial
\begin{itemize}
\item Sample: $n = 50$, $\bar{x} = 185$ mg/dL, $s = 30$ mg/dL
\item 95\% CI: $(176.5, 193.5)$ mg/dL
\end{itemize}

\vspace{0.5cm}
\textbf{We can immediately conclude:}
\begin{itemize}
\item $\mu = 200$? Not in CI, reject $H_0$ ($p < 0.05$)
\item $\mu = 190$? In CI, fail to reject $H_0$ ($p > 0.05$)
\item $\mu = 180$? In CI, fail to reject $H_0$ ($p > 0.05$)
\item $\mu = 170$? Not in CI, reject $H_0$ ($p < 0.05$)
\end{itemize}

\vspace{0.5cm}
\textbf{Interpretation:} Any value inside the CI is ``plausible'' given the data

Any value outside is ``implausible'' at the $\alpha = 0.05$ level
\end{frame}

\section{Practical Considerations}

\begin{frame}
\frametitle{Statistical vs Clinical Significance}
\textbf{Statistical significance:} $p < \alpha$

\textbf{Clinical significance:} Meaningful in practice

\vspace{0.5cm}
\textbf{These are NOT the same!}

\vspace{0.5cm}
\textbf{Example 1: Large sample, small effect}
\begin{itemize}
\item Drug lowers BP by 2 mmHg on average
\item With $n = 10000$: $p < 0.001$ (highly significant)
\item But 2 mmHg reduction may not be clinically meaningful
\end{itemize}

\vspace{0.5cm}
\textbf{Example 2: Small sample, large effect}
\begin{itemize}
\item Drug lowers BP by 15 mmHg on average
\item With $n = 10$: $p = 0.08$ (not significant at $\alpha = 0.05$)
\item But 15 mmHg reduction is clinically important!
\item Need larger sample to demonstrate statistical significance
\end{itemize}

\vspace{0.5cm}
\textbf{Always consider:} Effect size, confidence interval width, clinical context
\end{frame}

\begin{frame}
\frametitle{Sample Size and Hypothesis Testing}
\textbf{Effect of sample size on statistical tests:}

\vspace{0.5cm}
\textbf{Larger samples ($n \uparrow$):}
\begin{itemize}
\item Smaller standard error: $\text{SE} = s/\sqrt{n}$
\item Larger test statistics: $t = \frac{\bar{x} - \mu_0}{\text{SE}}$
\item Smaller p-values (easier to reject $H_0$)
\item Higher power to detect small effects
\item Narrower confidence intervals
\end{itemize}

\vspace{0.5cm}
\textbf{Implication:}
\begin{itemize}
\item With large enough $n$, even tiny effects become ``significant''
\item Always report:
  \begin{itemize}
  \item Effect size (e.g., difference in means)
  \item Confidence interval
  \item Clinical/practical importance
  \end{itemize}
\item Don't just report ``$p < 0.05$''!
\end{itemize}
\end{frame}

\begin{frame}
\frametitle{Common Mistakes in Hypothesis Testing}
\textbf{1. Confusing ``fail to reject'' with ``accept''}
\begin{itemize}
\item $\times$ ``We accept the null hypothesis''
\item \checkmark ``We fail to reject the null hypothesis''
\item Absence of evidence $\neq$ evidence of absence
\end{itemize}

\vspace{0.3cm}
\textbf{2. Choosing one-sided vs two-sided after seeing data}
\begin{itemize}
\item Must decide before collecting data based on research question
\end{itemize}

\vspace{0.3cm}
\textbf{3. Changing $\alpha$ after seeing p-value}
\begin{itemize}
\item ``p = 0.06, let's use $\alpha = 0.10$'' is p-hacking!
\end{itemize}

\vspace{0.3cm}
\textbf{4. Interpreting p-value as $\Pr(H_0 \text{ true})$}
\begin{itemize}
\item P-value is $\Pr(\text{data} \mid H_0)$, not $\Pr(H_0 \mid \text{data})$
\end{itemize}

\vspace{0.3cm}
\textbf{5. Ignoring assumptions}
\begin{itemize}
\item t-test assumes approximate normality (or large $n$)
\item Check with histogram, Q-Q plot, or normality tests
\end{itemize}

\vspace{0.3cm}
\textbf{6. Confusing statistical and practical significance}
\begin{itemize}
\item Always consider effect size and clinical importance
\end{itemize}
\end{frame}

\begin{frame}
\frametitle{Checking Assumptions}
\textbf{One-sample t-test assumes:}
\begin{enumerate}
\item Random sample from population
\item Observations are independent
\item Population is approximately normal (or $n$ large)
\end{enumerate}

\vspace{0.5cm}
\textbf{How to check normality:}

\vspace{0.3cm}
\textbf{1. Visual methods:}
\begin{itemize}
\item Histogram: Should be roughly bell-shaped
\item Q-Q plot: Points should follow straight line
\item Boxplot: Should be roughly symmetric
\end{itemize}

\vspace{0.3cm}
\textbf{2. Formal tests:}
\begin{itemize}
\item Shapiro-Wilk test: Tests $H_0$: data is normal
\item Anderson-Darling test
\item But: With large $n$, minor deviations become ``significant''
\end{itemize}

\vspace{0.3cm}
\textbf{3. Central Limit Theorem:}
\begin{itemize}
\item For $n \geq 30$, t-test is robust to non-normality
\item Exception: Extreme outliers or strong skewness
\end{itemize}
\end{frame}

\section{R Examples and Practice}

\begin{frame}[fragile]
\frametitle{One-Sample t-Test in R}
\textbf{Example: Body temperature data}

\vspace{0.2cm}
\small
\begin{verbatim}
# Body temperature data
temps <- c(98.0, 98.2, 98.4, 98.6, 98.0, 97.8, 98.8,
           98.0, 98.4, 98.6, 98.5, 98.3)

# Descriptive statistics
mean(temps)        # 98.30
sd(temps)          # 0.311
length(temps)      # 12

# Two-sided test: H0: mu = 98.6 vs HA: mu != 98.6
t.test(temps, mu = 98.6, alternative = "two.sided")

# Output includes:
#   t-statistic
#   degrees of freedom
#   p-value
#   95% confidence interval
#   sample mean
\end{verbatim}
\end{frame}

\begin{frame}[fragile]
\frametitle{One-Sided t-Test in R}
\small
\begin{verbatim}
# Blood pressure data (after treatment)
bp <- c(138, 142, 135, 140, 132, 128, 145, 136,
        139, 133, 141, 137, 134, 130, 143)

# One-sided test: H0: mu >= 140 vs HA: mu < 140
result <- t.test(bp, mu = 140, alternative = "less")

# Extract components
result$statistic    # t-statistic
result$parameter    # degrees of freedom
result$p.value      # p-value
result$conf.int     # one-sided confidence bound

# Manual calculation
xbar <- mean(bp)
s <- sd(bp)
n <- length(bp)
t_stat <- (xbar - 140) / (s / sqrt(n))
p_value <- pt(t_stat, df = n-1)

cat("Sample mean:", round(xbar, 2), "\n")
cat("t-statistic:", round(t_stat, 3), "\n")
cat("p-value:", round(p_value, 4), "\n")
\end{verbatim}
\end{frame}

\begin{frame}[fragile]
\frametitle{Checking Normality in R}
\small
\begin{verbatim}
# Generate some data
set.seed(524)
data <- rnorm(30, mean = 100, sd = 15)

# Visual checks
par(mfrow = c(1, 3))

# 1. Histogram
hist(data, breaks = 10, main = "Histogram",
     xlab = "Value", col = "lightblue")

# 2. Q-Q plot
qqnorm(data, main = "Q-Q Plot")
qqline(data, col = "red", lwd = 2)

# 3. Boxplot
boxplot(data, main = "Boxplot", ylab = "Value",
        col = "lightgreen")

par(mfrow = c(1, 1))

# Shapiro-Wilk test
shapiro.test(data)
# p > 0.05: fail to reject normality (data looks normal)
\end{verbatim}
\end{frame}

\begin{frame}[fragile]
\frametitle{Power and Sample Size in R}
\small
\begin{verbatim}
# Calculate power for one-sample t-test
# Scenario: detect difference of 5 units, SD = 10

library(pwr)  # May need to install: install.packages("pwr")

# Power for n = 30
pwr.t.test(n = 30,
           d = 5/10,        # effect size (delta/sigma)
           sig.level = 0.05,
           type = "one.sample",
           alternative = "two.sided")
# Power approx 0.70

# Sample size needed for 80% power
pwr.t.test(d = 5/10,
           power = 0.80,
           sig.level = 0.05,
           type = "one.sample",
           alternative = "two.sided")
# n approx 34
\end{verbatim}
\end{frame}

\begin{frame}[fragile]
\frametitle{Complete Hypothesis Test Example in R}
\footnotesize
\begin{verbatim}
# Research question: Is mean cholesterol > 200 mg/dL?
chol <- c(215, 198, 203, 225, 192, 210, 208, 195,
          218, 202, 212, 205, 220, 198, 207, 213,
          201, 209, 217, 204)

# Step 1: Descriptive statistics
cat("Sample size:", length(chol), "\n")
cat("Sample mean:", round(mean(chol), 2), "\n")
cat("Sample SD:", round(sd(chol), 2), "\n\n")

# Step 2: Check assumptions
shapiro.test(chol)  # Test normality

# Step 3: Conduct test (one-sided, upper tail)
result <- t.test(chol, mu = 200,
                 alternative = "greater",
                 conf.level = 0.95)

print(result)

# Step 4: Interpret
if(result$p.value < 0.05) {
  cat("\nReject H0. Evidence that mean cholesterol > 200.\n")
} else {
  cat("\nFail to reject H0. Insufficient evidence.\n")
}

# Step 5: Effect size
effect_size <- (mean(chol) - 200) / sd(chol)
cat("Cohen's d =", round(effect_size, 3), "\n")
\end{verbatim}
\end{frame}

\section{AI as a Learning Tool for Hypothesis Testing}

\AIcolors
\begin{frame}
\frametitle{AI for Hypothesis Testing}
\textbf{How AI can help this week:}
\begin{itemize}
\item \textbf{Conceptual clarification:} Understanding p-values, Type I/II errors
\item \textbf{Test selection:} Choosing appropriate test for your data
\item \textbf{R code generation:} Conducting t-tests, checking assumptions
\item \textbf{Interpretation:} Understanding output from \texttt{t.test()}
\item \textbf{Power analysis:} Sample size calculations
\item \textbf{Visualization:} Creating plots to check assumptions
\end{itemize}

\vspace{0.3cm}
\textbf{Example prompts:}
\begin{itemize}
\item ``Explain the difference between Type I and Type II errors with a medical example''
\item ``Show R code to conduct one-sided t-test testing if mean > 100''
\item ``How do I interpret a p-value of 0.03?''
\item ``Create Q-Q plot in R to check normality assumption''
\end{itemize}
\end{frame}

\begin{frame}[fragile]
\frametitle{AI Demo: Understanding P-values}
\textbf{Prompt:} ``Explain what a p-value of 0.02 means in plain language''

\vspace{0.2cm}
\textbf{Claude AI Response:}
\small
\begin{verbatim}
A p-value of 0.02 means:

"If the null hypothesis were actually true, we would
expect to see a result as extreme as (or more extreme
than) what we observed only 2% of the time by random
chance alone."

In practical terms:
- This is considered statistically significant
  (p < 0.05)
- It provides evidence against the null hypothesis
- There's only a 2% chance of getting this result if
  the null hypothesis is true
- We would typically reject the null hypothesis

Important: This does NOT mean:
- There's a 2% chance the null hypothesis is true
- There's a 98% chance our result is "real"
- The effect is necessarily large or important

Always consider the effect size and clinical
significance in addition to the p-value.
\end{verbatim}
\end{frame}

\begin{frame}[fragile]
\frametitle{AI Demo: Test Selection}
\textbf{Prompt:} ``I have data from 20 patients before and after treatment. What test should I use?''

\vspace{0.2cm}
\textbf{Claude AI Response:}
\tiny
\begin{verbatim}
For before-and-after data on the SAME patients, you should use a
PAIRED t-test (not a one-sample t-test).

Here's why:
- You have two measurements per patient (before and after)
- The measurements are dependent (same person)
- You're interested in the change within individuals

R code:
before <- c(...)  # before measurements
after <- c(...)   # after measurements

# Method 1: Paired t-test directly
t.test(after, before, paired = TRUE)

# Method 2: Create difference variable
difference <- after - before
t.test(difference, mu = 0)  # Test if mean difference != 0

Both methods give identical results. The paired t-test is more
powerful than comparing groups independently because it accounts
for individual variation.
\end{verbatim}
\end{frame}

\begin{frame}[fragile]
\frametitle{AI Demo: Checking Assumptions}
\textbf{Prompt:} ``How do I check if my data is normal enough for a t-test?''

\vspace{0.2cm}
\textbf{Claude AI Response:}
\tiny
\begin{verbatim}
# Check normality for t-test

data <- c(...)  # your data

# Visual checks (most important!)
par(mfrow = c(2, 2))

# 1. Histogram
hist(data, breaks = 10, main = "Histogram", prob = TRUE)
curve(dnorm(x, mean = mean(data), sd = sd(data)),
      add = TRUE, col = "red", lwd = 2)

# 2. Q-Q plot (most useful!)
qqnorm(data)
qqline(data, col = "red")
# Points should roughly follow the line

# 3. Boxplot
boxplot(data, main = "Boxplot")
# Should be roughly symmetric, check for outliers

# 4. Density plot
plot(density(data), main = "Density Plot")

# Formal test
shapiro.test(data)
# p > 0.05 suggests normality
# But with large n, minor deviations become "significant"

# Rule of thumb: t-test is robust for n >= 30 unless
# extreme outliers or strong skewness
\end{verbatim}
\end{frame}

\begin{frame}
\frametitle{AI Best Practices for Hypothesis Testing}
\begin{columns}
\begin{column}{0.5\textwidth}
\textbf{DO:}
\begin{itemize}
\item Ask AI to explain concepts
\item Request code examples
\item Verify calculations independently
\item Check assumptions visually
\item Ask for interpretation help
\item Learn from AI explanations
\end{itemize}
\end{column}

\begin{column}{0.5\textwidth}
\textbf{DON'T:}
\begin{itemize}
\item Blindly trust AI p-values
\item Skip checking assumptions
\item Use AI without understanding
\item Choose test based on desired result
\item Ignore effect sizes
\item Forget clinical context
\end{itemize}
\end{column}
\end{columns}

\vspace{0.3cm}
\textbf{Critical thinking points:}
\begin{enumerate}
\item Does the test AI suggested match your research question?
\item Are assumptions reasonable for your data?
\item Does the interpretation make practical sense?
\item Is statistical significance the same as clinical importance?
\end{enumerate}
\end{frame}

\defaultcolors

\section{Summary}

\begin{frame}
\frametitle{Summary: Hypothesis Testing Framework}
\textbf{Complete hypothesis testing procedure:}

\vspace{0.3cm}
\begin{enumerate}
\item \textbf{State hypotheses:} $H_0$ vs $H_A$ (before collecting data!)
\item \textbf{Choose significance level:} Usually $\alpha = 0.05$
\item \textbf{Check assumptions:} Normality, independence, random sampling
\item \textbf{Calculate test statistic:} $t = \frac{\bar{x} - \mu_0}{s/\sqrt{n}}$
\item \textbf{Find p-value:} From t-distribution with $n-1$ df
\item \textbf{Make decision:} Reject $H_0$ if $p < \alpha$
\item \textbf{State conclusion:} In context of research question
\item \textbf{Report effect size:} Difference, confidence interval, practical importance
\end{enumerate}

\vspace{0.5cm}
\textbf{Remember:}
\begin{itemize}
\item Statistical significance $\neq$ practical importance
\item Fail to reject $\neq$ accept
\item Always consider Type I and Type II error consequences
\item Report confidence intervals alongside p-values
\end{itemize}
\end{frame}

\begin{frame}
\frametitle{Key Concepts to Remember}
\textbf{Hypotheses:}
\begin{itemize}
\item $H_0$: null hypothesis (status quo, no effect)
\item $H_A$: alternative hypothesis (what we want to show)
\item Two-sided vs one-sided tests
\end{itemize}

\vspace{0.3cm}
\textbf{Decision-making:}
\begin{itemize}
\item P-value: probability of data or more extreme, given $H_0$ true
\item $\alpha$: significance level (Type I error rate)
\item Reject $H_0$ if $p < \alpha$
\end{itemize}

\vspace{0.3cm}
\textbf{Errors:}
\begin{itemize}
\item Type I ($\alpha$): False positive (reject true $H_0$)
\item Type II ($\beta$): False negative (fail to reject false $H_0$)
\item Power = $1 - \beta$: Probability of detecting true effect
\end{itemize}

\vspace{0.3cm}
\textbf{Tests covered:}
\begin{itemize}
\item One-sample t-test: $t = \frac{\bar{x} - \mu_0}{s/\sqrt{n}}$
\item One-sample z-test: $Z = \frac{\bar{x} - \mu_0}{\sigma/\sqrt{n}}$ (rare)
\end{itemize}
\end{frame}

\begin{frame}
\frametitle{Next Week Preview}
\textbf{Week 6: Two-Sample Tests}

\vspace{0.5cm}
We'll extend hypothesis testing to compare two groups:
\begin{itemize}
\item Two-sample t-test (independent groups)
\item Paired t-test (dependent samples)
\item Confidence intervals for difference in means
\item Assumptions and diagnostics
\item Effect sizes (Cohen's d)
\end{itemize}

\vspace{0.5cm}
\textbf{Examples:}
\begin{itemize}
\item Treatment vs control groups
\item Before vs after measurements
\item Drug A vs Drug B comparison
\end{itemize}

\vspace{0.5cm}
\textbf{Reading:} Chapter 11
\end{frame}

\begin{frame}
\frametitle{Questions?}
\begin{center}
\Large{Thank you!}

\vspace{1cm}

\textbf{Email:} john.molitor@oregonstate.edu\\
\textbf{Office:} 153 Milam
\end{center}
\end{frame}

\end{document}
